\section{Threats to Validity}
\label{sec:threats}

\textbf{External validity.} All measurements come from a single laptop running Spark in local
mode. There is no network shuffle and no cluster scheduling, so absolute times do not carry
over to production clusters. The break-even reasoning and the relative effects of the actions
are what we expect to carry over, and the framework itself consumes measured or estimated
costs rather than fixed constants. The datasets are small (at most 28\,MB on disk and 35\,MB
shuffled); with larger data the best partition count will grow, and our specific optima should
not be reused as settings. The data are synthetic, and the profiled workloads use a single
data seed; to limit this threat, the validation workloads use different seeds and a volume not
seen during profiling.

\textbf{Platform effects.} The experiments run on Windows with NTFS, where creating many small
files is comparatively expensive. This amplifies the overhead of Spark's bypass-merge shuffle
writer, which Spark uses when a shuffle has no map-side aggregation and at most 200 reduce
partitions
(\texttt{spark.shuffle.sort.bypassMergeThreshold} \cite{spark2024config}) and which writes one
file per reduce partition for every map task. It contributes to
both the size of the A1 effect and the mechanism behind A2. On Linux, HDFS or object storage we
expect the effects to be smaller, but we have no reason to expect them to reverse. The hybrid
processor (performance and efficiency cores) caused occasional slow executions at 1--2
partitions, raising the coefficient of variation above 0.20 in 5 of the 32 A1 cells. We report
these cells as measured, and a sensitivity analysis that removes the slowest execution of each
flagged cell leaves the conclusions unchanged.

\textbf{Internal validity.} Disabling AQE and broadcast joins was necessary to control the
partition count and to force a shuffle, but it departs from Spark's defaults; AQE in particular
would reduce the benefit of manual partition tuning. The A2 benefit is partly an interaction
with the static setting of 200 shuffle partitions: fragmented input produces more map tasks,
and each of them writes 200 shuffle files. Because the design measures each action alone, the
joint effect of A1 and A2 is not quantified. In A2 the warm-up executions leave input files in
the operating system's cache, so the measured benefit of compaction is a lower bound for cold
reads from remote storage. We mitigated systematic drift with session warm-ups, a discarded
warm-up before every measured execution, randomized block order with fixed seeds, and
pre-flight checks on free memory and power supply.

\textbf{Construct validity.} Costs are estimated from published on-demand prices and are never
billed. Proxy P1 assumes the whole machine is paid for while the job runs, whereas P2 charges
only the core time that tasks use; neither includes storage, data-transfer or licence costs.
We therefore report both proxies and vary the price by $\pm 50\,\%$. The number of future
executions $N$ is a model parameter and is not physically executed, so the model assumes the
measured saving recurs unchanged. Charging the A1 search entirely as overhead overstates its
one-time cost; the rule-based variant bounds it from below.

\textbf{Conclusion validity.} With five blocks per cell, rank-based significance tests have
little power, so we rely on medians, interquartile ranges, the consistency of the paired
savings across blocks, and bootstrap confidence intervals. Differences smaller than the noise
floor are treated as ties rather than as effects.
